\documentclass{article}
\usepackage{amsmath}
\usepackage{amssymb}
\usepackage{amsthm}

\theoremstyle{definition}
\newtheorem{definisi}{Definisi}
\newtheorem{teorema}{Teorema}
\newtheorem*{bukti}{Bukti}
\newtheorem*{sangkal}{Sangkal}

\begin{document}
\section*{POV Himpunan}
\begin{definisi}[Himpunan Indeks dan Skalar]
  Misalkan $N = \{1, 2, \dots, n\}$ adalah himpunan indeks koordinat matriks. Misalkan $S \subset \mathbb{R}$ adalah himpunan skalar berhingga yang menyusun Latin square $A$ dan $B$, sedemikian sehingga elemen matriks $A_{ij}, B_{ij} \in S$ untuk setiap $i, j \in N$.
\end{definisi}

\begin{definisi}[Dekomposisi Permutasi Matriks]
  Untuk setiap nilai skalar $x \in S$, letak elemen $x$ pada matriks Latin square $A$ merepresentasikan sebuah permutasi unik. Definisikan fungsi permutasi $P_x : N \to N$ sedemikian sehingga:
  $$A_{i, P_x(i)} = x \quad \text{untuk setiap } i \in N$$
  Secara serupa, definisikan fungsi permutasi $Q_y : N \to N$ untuk matriks $B$ sedemikian sehingga:
  $$B_{i, Q_y(i)} = y \quad \text{untuk setiap } i \in N$$
\end{definisi}

\begin{definisi}[Representasi Koordinat Permutasi]
  Untuk setiap fungsi permutasi $\pi : N \to N$, definisikan representasi koordinatnya, $\Gamma(\pi)$, sebagai himpunan pasangan berurutan:
  $$\Gamma(\pi) = \{(i, \pi(i)) \mid i \in N\}$$
\end{definisi}

\begin{definisi}[Himpunan Tingkat Atas Kumulatif]
  Untuk setiap nilai batas observasi $v \in \mathbb{R}$, definisikan $\mathcal{U}_{\ge v}^{AB}$ sebagai gabungan dari seluruh representasi koordinat komposisi permutasi dari perkalian $A \otimes B$ yang menghasilkan bobot skalar $x + y \ge v$:
  $$\mathcal{U}_{\ge v}^{AB} = \bigcup_{\substack{x, y \in S \\ x+y \ge v}} \Gamma(Q_y \circ P_x)$$
  Dan secara simetris untuk arah $B \otimes A$:
  $$\mathcal{U}_{\ge v}^{BA} = \bigcup_{\substack{x, y \in S \\ x+y \ge v}} \Gamma(P_x \circ Q_y)$$
\end{definisi}

\begin{teorema}[Kesetaraan Himpunan Tingkat]
  Dua matriks Latin square $A$ dan $B$ saling komutatif terhadap operasi perkalian max-plus ($A \otimes B = B \otimes A$) \textbf{jika dan hanya jika} untuk setiap nilai batas observasi $v$, himpunan tingkat atas dari kedua arah perkalian tersebut ekuivalen:
  $$\mathcal{U}_{\ge v}^{AB} = \mathcal{U}_{\ge v}^{BA}$$
\end{teorema}

\begin{bukti}
  Misalkan $C = A \otimes B$ dan $D = B \otimes A$.
  Berdasarkan definisi perkalian max-plus, elemen pada matriks $C$ adalah:
  $$C_{ij} = \max_{k \in N} (A_{ik} + B_{kj})$$

  Kombinasi nilai $A_{ik} = x$ dan $B_{kj} = y$ terjadi jika dan hanya jika $k = P_x(i)$ dan $j = Q_y(k)$.
  Maka nilai bobot $C_{ij} \ge v$ dapat tercapai jika dan hanya jika terdapat skalar $x, y \in S$ dengan $x+y \ge v$ sedemikian sehingga $j = (Q_y \circ P_x)(i)$.
  Secara ekuivalen dalam teori himpunan, hal ini berarti $C_{ij} \ge v \iff (i, j) \in \mathcal{U}_{\ge v}^{AB}$.

  Dengan logika yang simetris dari $D_{ij} = \max_{k} (B_{ik} + A_{kj})$, elemen pada matriks $D$ memenuhi $D_{ij} \ge v \iff (i, j) \in \mathcal{U}_{\ge v}^{BA}$.
  \begin{itemize}
    \item \textbf{($\Rightarrow$)} Asumsikan $A \otimes B = B \otimes A$, yang berarti $C_{ij} = D_{ij}$ untuk setiap $i, j \in N$.
          Ambil sembarang nilai batas $v \in \mathbb{R}$.
          Jika koordinat $(i, j) \in \mathcal{U}_{\ge v}^{AB}$, maka $C_{ij} \ge v$.
          Karena matriks komutatif ($C_{ij} = D_{ij}$), maka $D_{ij} \ge v$, yang secara definitif mengimplikasikan $(i, j) \in \mathcal{U}_{\ge v}^{BA}$.
          Sebaliknya, jika $(i, j) \in \mathcal{U}_{\ge v}^{BA}$, maka $D_{ij} \ge v \implies C_{ij} \ge v \implies (i, j) \in \mathcal{U}_{\ge v}^{AB}$.
          Oleh karena itu, terbukti bahwa $\mathcal{U}_{\ge v}^{AB} = \mathcal{U}_{\ge v}^{BA}$.

    \item \textbf{($\Leftarrow$)} Asumsikan kesetaraan himpunan tingkat berlaku, yaitu $\mathcal{U}_{\ge v}^{AB} = \mathcal{U}_{\ge v}^{BA}$ untuk setiap $v \in \mathbb{R}$.
          Kita akan membuktikannya menggunakan kontradiksi. Andaikan $A \otimes B \neq B \otimes A$, maka haruslah terdapat setidaknya satu pasang koordinat $(i, j)$ di mana $C_{ij} \neq D_{ij}$.
          Tanpa mengurangi keumuman, asumsikan $C_{ij} > D_{ij}$.

          Misalkan $v = C_{ij}$. Karena nilai maksimum di koordinat $(i, j)$ pada matriks $C$ adalah tepat sebesar $v$, maka $(i, j) \in \mathcal{U}_{\ge v}^{AB}$.
          Namun di sisi lain, karena $D_{ij} < v$, maka rute pada koordinat $(i, j)$ tidak mungkin mencapai bobot sebesar $v$ pada matriks $D$. Akibatnya, $(i, j) \notin \mathcal{U}_{\ge v}^{BA}$.
          Hal ini menghasilkan kontradiksi langsung dengan asumsi awal kita bahwa $\mathcal{U}_{\ge v}^{AB} = \mathcal{U}_{\ge v}^{BA}$.
          Maka, pengandaian tersebut salah. Haruslah $C_{ij} = D_{ij}$ untuk semua $i, j \in N$.
          Terbukti bahwa $A \otimes B = B \otimes A$.
  \end{itemize}

  \section{False Theorem}
  \begin{teorema}
    Jika dua buah latin square hasil operasi $\otimes$ nya menghasilkan Latin square juga maka dua buah latin square itu komutatif
  \end{teorema}
  \begin{sangkal}
    Misalkan $A$ dan $B$ adalah dua buah matriks Latin square berukuran $3 \times 3$ berikut:

    $$A = \begin{pmatrix} 1 & 3 & 2 \\ 3 & 2 & 1 \\ 2 & 1 & 3 \end{pmatrix}, \quad B = \begin{pmatrix} 1 & 2 & 3 \\ 2 & 3 & 1 \\ 3 & 1 & 2 \end{pmatrix}$$

    $$A \otimes B = \begin{pmatrix} 5 & 6 & 4 \\ 4 & 5 & 6 \\ 6 & 4 & 5 \end{pmatrix}$$

    $$B \otimes A = \begin{pmatrix} 5 & 4 & 6 \\ 6 & 5 & 4 \\ 4 & 6 & 5 \end{pmatrix}$$

  \end{sangkal}
\end{bukti}

\end{document}